%!TEX root = thesis.tex

\chapter{Conclusion}
\label{ch:conclusion}

\section{Summary}
\label{sec:concl:summary}

This Praktikum report presented the design and implementation of a \ac{VL} polyp segmentation framework that integrates structured clinical text descriptions with visual features to improve both segmentation quality and uncertainty estimation. The key components of the system include:

\begin{itemize}
	\item A \textbf{\ac{ResNet}-34 vision encoder} pretrained on ImageNet, providing multi-scale feature extraction across five resolution levels.
	
	\item A \textbf{frozen \ac{PubMedBERT} text encoder} that processes structured clinical captions derived from the Paris Classification system, capturing morphology, size, location, and surface characteristics.
	
	\item A \textbf{cross-attention fusion module} with 8 attention heads that enables fine-grained vision-language interaction, allowing spatial visual features to selectively attend to relevant clinical text information.
	
	\item A \textbf{U-Net-style decoder with \ac{MC} Dropout} ($p = 0.1$, $T = 10$ samples) for simultaneous segmentation and uncertainty quantification.
	
	\item \textbf{Three novel evaluation metrics} (\ac{URS}, \ac{SAS}, \ac{AC-ECE}) specifically designed to quantify the impact of language guidance on uncertainty quality.
\end{itemize}

The framework was designed with modularity and reproducibility as guiding principles. The configuration system supports six experimental presets that enable systematic ablation across clinical text attributes.


\section{Findings}
\label{sec:concl:findings}

The experimental framework has been fully implemented and validated. Detailed quantitative conclusions will be drawn upon completion of all experimental runs. The ablation study across six configurations (vision-only, full VL, and four single-attribute variants) will provide insights into:

\begin{enumerate}
	\item The overall effectiveness of language guidance for polyp segmentation.
	\item The relative importance of individual clinical attributes (shape, size, location, pathology).
	\item The impact of text integration on uncertainty calibration across different polyp types.
	\item Cross-dataset generalization to Kvasir-SEG and CVC-ClinicDB.
\end{enumerate}


\section{Limitations}
\label{sec:concl:limitations}

Several limitations of the current approach should be noted:

\begin{itemize}
	\item \textbf{Text availability at inference}: The framework assumes access to clinical metadata for generating text descriptions at test time. In real deployment, this metadata may not be available, requiring either a default caption strategy or a separate metadata prediction module.
	
	\item \textbf{Frozen text encoder}: While freezing PubMedBERT prevents overfitting, it also prevents the text representations from adapting to the specific visual characteristics of the polyp segmentation task.
	
	\item \textbf{Dataset size}: The SUN dataset, while substantial (49{,}136 frames), is derived from only 100 unique polyps. The effective diversity is limited by the video-based nature of the data.
	
	\item \textbf{MC Dropout approximation}: MC Dropout provides a computationally efficient but approximate uncertainty estimate. More sophisticated methods (e.g., deep ensembles) may yield higher-quality uncertainties at increased computational cost.
\end{itemize}


\section{Future Work}
\label{sec:concl:future}

Several directions for future work emerge from this project:

\begin{itemize}
	\item \textbf{Learnable text encoder}: Fine-tuning PubMedBERT (or a portion of it) with appropriate regularization and potentially with a larger medical image-text dataset.
	
	\item \textbf{Automatic caption generation}: Training a separate model to predict clinical attributes directly from the input image, removing the dependency on external metadata.
	
	\item \textbf{Stronger vision backbones}: Replacing ResNet-34 with more powerful encoders such as ResNet-101 or Vision Transformers (ViT), and investigating the interaction between backbone capacity and language benefit.
	
	\item \textbf{Multi-scale fusion}: Applying cross-attention at multiple encoder levels rather than only at the bottleneck, potentially capturing attribute-relevant features at different spatial scales.
	
	\item \textbf{Real-time inference}: Optimizing the pipeline for real-time deployment during colonoscopy, potentially through model distillation, pruning, or efficient attention mechanisms.
	
	\item \textbf{Extended evaluation}: Testing on additional polyp segmentation benchmarks (e.g., ETIS-Larib, ColonDB) and evaluating with boundary-specific metrics such as boundary IoU.
\end{itemize}
